% Fragmento público generado desde propietarios íntegros.
% Residencia: 52.11.1.
% Fuente: ../incoming/radion_monografia_20260827/manuscrito/sections/12_aritmetica_borde.tex:1-159
\noindent La frontière nonadique conserve le défaut somme--produit, la vacance et la retenue de complétion.\par\medskip

\subsubsection{Radical \(3\! -\! 6\! -\! 9\) du bord}

\paragraph{Génération et nilpotence du radical nonadique}

Dans l'anneau \(\Z/9\Z\), l'idéal
\begin{equation}
\mathfrak m=(3)=\{0,3,6\}
\label{eq:radical-369}
\end{equation}
satisfait
\begin{equation}
\mathfrak m^2=(9)=0.
\label{eq:m-square-zero}
\end{equation}
C'est le secteur nilpotent du bord nonadique. L'orbite visible de ses trois
marques se publie comme
\[
369\longrightarrow693\longrightarrow936\longrightarrow369.
\]
Le retour n'implique pas que tous les états soient égaux : la permutation
visible se ferme tandis que les cocycles de chemin peuvent avancer.

La double convention
\[
9\equiv0\pmod9,
\qquad
\operatorname{dr}_9(9)=9
\]
exprime deux lecteurs : zéro résiduel algébrique et clôture visible de la racine
numérique. C'est pourquoi un neuf peut absorber une torsion de bord sans devenir
un zéro spécifique de la fonction zêta.

\paragraph{Défaut somme--produit \(459-405=54\)}

Les feuillets APP produisent des totaux complémentaires
\[
S_+=405,\qquad S_\times=459,\qquad S_\times-S_+=54.
\]
Le \(54\) est un défaut exact des tables, et réapparaît comme exposant de
contraction par l'identité \(9\cdot6=54\). Cependant, chaque occurrence
nécessite son application. La distance focale \(d_1=0.544545\ldots\) n'est pas égale à
\(0.54\), et l'analyse des invariants n'a pas produit de flèche qui
l'identifie au défaut \(54\).

La puissance
\[
9^9=387,420,489
\]
génère une échelle où
\begin{equation}
\operatorname{ord}_{10}\bigl(9^{9^9}\bigr)=369,693,099,
\qquad
\#\,\text{chiffres}=369,693,100.
\label{eq:369693}
\end{equation}
Le bloc \(369|693\) apparaît dans le couple ordre--longueur ; on n'affirme pas qu'il soit
le préfixe décimal de la puissance.

\subsubsection{Théorie arithmétique des vacances décimales}

\paragraph{Pente asymptotique de vacance}

La séparation entre la décennie et la base nonadique est
\begin{equation}
\delta_{\mathrm{vac}}=\log_{10}\frac{10}{9}.
\label{eq:delta-vac}
\end{equation}
Le compteur cumulé
\begin{equation}
V_9(n)=\left\lceil n\delta_{\mathrm{vac}}\right\rceil
\label{eq:vacancy-count}
\end{equation}
enregistre le nombre de positions de correction nécessaires à la publication décimale
d'une évolution nonadique. Ses sauts forment un mot binaire
\[
z_n=V_9(n)-V_9(n-1)\in\{0,1\}.
\]
Les lacunes se répartissent avec des longueurs \(21/22\), les retours avec des
longueurs \(6/7\), et le déterminant combiné est \(15\). Ces nombres
proviennent du lecteur de vacance, non du noyau de Barbero, même si le \(15\)
coïncide avec \(\binom62\).

\paragraph{Polarisation et courant arithmétique du réseau nonadique}

On définit
\begin{equation}
P_N=\sum_{n=1}^N(z_n-\delta_{\mathrm{vac}}),
\qquad
J_N=P_N-P_{N-1}=z_N-\delta_{\mathrm{vac}}.
\label{eq:arith-current}
\end{equation}
La séquence équilibrée satisfait \(|P_N|<1\). Le \enquote{courant}
\(J_N\) est ici un courant arithmétique : il mesure la création et l'absorption de
vacances par rapport à la pente moyenne. Seul un transducteur dimensionnel
ultérieur peut le publier comme densité de courant électromagnétique.

L'intuition physique est puissante parce que la même forme algébrique sépare une
moyenne paire et une fluctuation orientée. La rigueur exige de conserver le type :
\[
J_N\in\R\quad\text{sans dimension},
\qquad
j^\mu\in\text{droite physique de courant}.
\]

\subsubsection{Complétion \(729\to1000\) et conservation de la retenue}

\paragraph{Décomposition \(1000=729+243+27+1\)}

La complétion
\[
1000=729+243+27+1
\]
conserve l'unité au moyen de quatre échelles ternaires. La couronne totale est
\[
1000-729=271,
\]
tandis que le dernier entier avant la clôture satisfait
\[
999=729+270.
\]
La différence entre \(270\) et \(271\) distingue frontière occupée et unité
complétante. Le même soin évite d'identifier une vacance à un zéro ou une
retenue à un résidu réel.

\paragraph{Résidence de la retenue du radion dans la couronne décimale}

Les opérandes du radion commencent, en blocs de trois chiffres, par
\[
\Phi_T:\quad000|027|455|906|837|565|\cdots,
\]
\[
D_E:\quad000|027|455|010|034|743|\cdots.
\]
La soustraction APP donne
\[
\epsone:\quad000|000|000|896|802|822|\cdots.
\]
La petitesse du résultat ne s'obtient pas en tronquant les termes inférieurs.
Elle s'obtient parce que deux publications distinctes partagent un long préfixe et que
l'opérateur \(\operatorname{SubAcarreo}_{1000}\) conserve les emprunts nécessaires
pour les soustraire exactement.

Le mot \emph{mémoire} a ici deux niveaux :
\begin{enumerate}
\item \(q_{\mathrm{mem}}\) enregistre la profondeur des retours nonadiques ;
\item \(c_i\) enregistre les emprunts locaux de la soustraction décimale.
\end{enumerate}
Le premier organise l'hélice ; le second publie la décimale. \(\epsR\) utilise
les deux parce que ses opérandes proviennent de cylindres en profondeur et que sa soustraction
s'effectue avec retenue, mais il n'est égal à aucun des compteurs.

\begin{thesisbox}
Le bord n'est pas un lieu où la théorie perd de l'information. C'est le lieu où
l'information change de représentation : phase qui revient, mémoire qui
s'élève, cylindre qui se contracte, vacance qui s'enregistre et retenue qui
transporte la différence.
\end{thesisbox}
